% ТЗ: Граф-аналитик — xelatex graph_analyzer_tz.tex
\documentclass[11pt, a4paper, oneside]{article}

\usepackage{fontspec}
\usepackage{polyglossia}
\setdefaultlanguage{russian}
\setotherlanguage{english}
\setmainfont{Liberation Serif}
\setsansfont{Liberation Sans}
\setmonofont[Scale=0.85]{DejaVu Sans Mono}

\usepackage[
  a4paper,
  top=22mm, bottom=22mm,
  left=28mm, right=22mm,
  headheight=14pt
]{geometry}

\usepackage[protrusion=false,expansion=false]{microtype}
\usepackage{parskip}
\usepackage{setspace}
\usepackage{booktabs}
\usepackage{array}
\usepackage{tabularx}
\usepackage{enumitem}
\usepackage{graphicx}
\usepackage{float}
\usepackage{caption}
\usepackage{amsmath}
\usepackage{amssymb}
\usepackage{listings}

\usepackage[dvipsnames, svgnames]{xcolor}
\definecolor{AccentBlue}{RGB}{25, 90, 160}
\definecolor{AccentTeal}{RGB}{0, 115, 115}
\definecolor{AccentOrange}{RGB}{190, 70, 0}
\definecolor{AccentRed}{RGB}{160, 20, 20}
\definecolor{AccentGreen}{RGB}{28, 130, 55}
\definecolor{LightGray}{RGB}{246, 248, 252}
\definecolor{MidGray}{RGB}{120, 125, 138}
\definecolor{DarkGray}{RGB}{35, 38, 48}
\definecolor{CodeBg}{RGB}{244, 246, 252}
\definecolor{CodeBorder}{RGB}{195, 202, 222}

\usepackage[
  unicode, colorlinks=true,
  linkcolor=AccentBlue,
  urlcolor=AccentTeal,
  pdftitle={ТЗ: Граф-аналитик процессов},
  pdfauthor={development\_system}
]{hyperref}

\lstset{
  basicstyle=\small\ttfamily\color{DarkGray},
  backgroundcolor=\color{CodeBg},
  frame=single, framerule=0.5pt, rulecolor=\color{CodeBorder},
  breaklines=true, showspaces=false, showstringspaces=false,
  numbers=left, numberstyle=\tiny\color{MidGray}, numbersep=8pt,
  xleftmargin=14pt, xrightmargin=4pt,
  aboveskip=8pt, belowskip=4pt, tabsize=2,
}



\usepackage{fancyhdr}
\pagestyle{fancy}
\fancyhf{}
\fancyhead[L]{\small\color{MidGray}\sffamily Граф-аналитик процессов}
\fancyhead[R]{\small\color{MidGray}\sffamily Техническое задание}
\fancyfoot[C]{\small\color{MidGray}\thepage}
\renewcommand{\headrulewidth}{0.4pt}
\renewcommand{\headrule}{\hbox to\headwidth{\color{AccentBlue}\leaders\hrule height 0.4pt\hfill}}
\renewcommand{\footrulewidth}{0pt}

\usepackage{titlesec}
\titleformat{\section}
  {\large\bfseries\sffamily\color{AccentBlue}}
  {\thesection.}{0.7em}{}
  [\vspace{1pt}{\color{AccentBlue}\rule{\linewidth}{0.5pt}}\vspace{3pt}]
\titleformat{\subsection}
  {\normalsize\bfseries\color{DarkGray}}
  {\thesubsection.}{0.6em}{}
\titleformat{\subsubsection}
  {\small\bfseries\color{AccentTeal}}
  {\thesubsubsection.}{0.5em}{}

\titlespacing*{\section}{0pt}{16pt}{6pt}
\titlespacing*{\subsection}{0pt}{10pt}{3pt}
\titlespacing*{\subsubsection}{0pt}{7pt}{2pt}

\usepackage{tcolorbox}
\tcbuselibrary{breakable, skins}

\newtcolorbox{infobox}[1][]{
  colback=LightGray, colframe=AccentBlue,
  colbacktitle=AccentBlue, coltitle=white,
  fonttitle=\small\bfseries\sffamily,
  boxrule=0.5pt, arc=3pt, breakable,
  top=4pt, bottom=4pt, left=6pt, right=6pt, #1}

\newtcolorbox{warnbox}[1][]{
  colback=Seashell, colframe=AccentOrange,
  colbacktitle=AccentOrange, coltitle=white,
  fonttitle=\small\bfseries\sffamily,
  boxrule=0.5pt, arc=3pt, breakable,
  top=4pt, bottom=4pt, left=6pt, right=6pt, #1}

\newtcolorbox{okbox}[1][]{
  colback=Honeydew, colframe=AccentGreen,
  colbacktitle=AccentGreen, coltitle=white,
  fonttitle=\small\bfseries\sffamily,
  boxrule=0.5pt, arc=3pt, breakable,
  top=4pt, bottom=4pt, left=6pt, right=6pt, #1}

\newtcolorbox{reqbox}[2][]{
  colback=LightGray, colframe=AccentBlue!40,
  fontupper=\small,
  title={\small\sffamily\color{AccentBlue}\textbf{#2}},
  boxrule=0.4pt, arc=2pt, breakable,
  top=3pt, bottom=3pt, left=6pt, right=6pt, #1}

\newcolumntype{L}{>{\raggedright\arraybackslash}X}
\renewcommand{\arraystretch}{1.3}
\setlength{\parskip}{5pt}
\setlength{\parindent}{0pt}
\setlist{nosep, leftmargin=1.4em, topsep=3pt}
\captionsetup{font=small, labelfont=bf}

% ════════════════════════════════════════════════════════════════
\begin{document}

% ── Титульная страница ───────────────────────────────────────────
\begin{titlepage}
\centering
\vspace*{18mm}

{\color{AccentBlue}\rule{\textwidth}{1.5pt}}
\vspace{8mm}

{\sffamily\fontsize{10}{12}\color{MidGray} ТЕХНИЧЕСКОЕ ЗАДАНИЕ}\\[6mm]

{\sffamily\fontsize{22}{28}\bfseries\color{AccentBlue}
  Граф-аналитик процессов
}\\[4mm]

{\sffamily\fontsize{13}{16}\color{DarkGray}
  Инструмент анализа бизнес-процессов:\\
  поиск крайних случаев, пробелов компенсации\\
  и неожиданных состояний
}

\vspace{10mm}
{\color{AccentBlue}\rule{\textwidth}{0.4pt}}
\vspace{10mm}

\begin{minipage}{0.85\textwidth}
\centering\itshape\color{DarkGray}\fontsize{11}{15}\selectfont
Человек описывает бизнес-процесс в виде Mermaid-графа
и передаёт ключевые Django-модели. Инструмент использует LLM
для анализа графа и возвращает структурированный отчёт:
неожиданные ситуации, крайние случаи, пробелы в компенсации
и недостижимые ветки — \textbf{до написания кода и тестов}.
\end{minipage}

\vfill

\begin{tabular}{ll}
  \textbf{Продукт:}       & Граф-аналитик (поток в development\_system) \\[3pt]
  \textbf{Версия ТЗ:}     & 1.0 \\[3pt]
  \textbf{Дата:}          & 14 марта 2026 \\[3pt]
  \textbf{Стек:}          & Python / Common Lisp, OpenRouter LLM \\[3pt]
  \textbf{Применение:}    & Анализ add\_tour и аналогичных процессов \\
\end{tabular}

\vspace{10mm}
{\color{AccentBlue}\rule{\textwidth}{0.4pt}}
{\small\color{MidGray}
  Mermaid · Django ORM · LLM · development\_system · Graph Analysis
}
\end{titlepage}

\tableofcontents
\newpage

% ════════════════════════════════════════════════════════════════
\section{Назначение}

\subsection{Контекст}

При разработке сложных бизнес-процессов (оркестрация задач, многофазные
транзакции, условные переходы) разработчик описывает логику как граф
состояний и переходов. Такой граф неизбежно содержит:

\begin{itemize}
  \item \textbf{Крайние случаи} — входные данные или состояния БД,
    которые граф не рассматривает явно
  \item \textbf{Пробелы компенсации} — ветки ошибок, не откатывающие
    часть уже выполненных шагов
  \item \textbf{Неожиданные состояния} — комбинации полей модели,
    нарушающие предположения графа
  \item \textbf{Недостижимые ветки} — узлы или ветки, до которых
    алгоритм никогда не доберётся при нормальной работе
\end{itemize}

Сейчас такой анализ выполняется вручную (code review, опыт команды)
и пропускает значительную часть проблем до попадания в продакшн.

\subsection{Цель инструмента}

\begin{infobox}[title={Цель}]
Принять граф процесса (Mermaid) и ключевые Django-модели,
проанализировать их с помощью LLM и вернуть структурированный
отчёт о слабых местах \textbf{до написания тестов и кода}.
\end{infobox}

Инструмент не заменяет code review и не гарантирует полноту анализа.
Он снижает стоимость обнаружения проблем, перемещая их из фазы тестирования
в фазу проектирования.

\subsection{Первичный прецедент использования}

Процесс \texttt{add\_tour} — добавление контейнеров в маршрут.
45 событий, 5 фаз, несколько точек компенсации, версионирование
контейнеров. Ключевая модель: \texttt{WasteSiteVisitedContainers}.

% ════════════════════════════════════════════════════════════════
\section{Входные данные}

\subsection{Граф процесса (обязательно)}

Формат: Mermaid \texttt{flowchart}. Требования к графу:

\begin{itemize}
  \item Узлы имеют идентификаторы (\texttt{A0}, \texttt{B1}, ...)
  \item Метки узлов описывают действие (например: \texttt{E12 Обновить статус контейнеров})
  \item Рёбра имеют условия на ветвлениях (\texttt{|"Да"|}, \texttt{|"Нет"|})
  \item Узлы ошибок явно помечены (\texttt{ERR0}, \texttt{ERR1} или аналог)
  \item Терминальные узлы обозначены (\texttt{OK}, \texttt{Финиш}, ...)
\end{itemize}

\begin{lstlisting}[language={}, caption={Пример входного графа (фрагмент)}]
flowchart TD
  A6["E6 Проверка зон РО"] --> A6C{"Ок?"}
  A6C -->|"Нет"| ERR0["Компенсация и ошибка"]
  A6C -->|"Да"| A7["E7 Проверка совместимости контейнеров с ТС"]
\end{lstlisting}

\subsection{Django-модели (обязательно)}

Передаются как Python-код: определения классов с полями и их типами.
Достаточно полей участвующих в процессе — не нужны все поля модели.

\begin{lstlisting}[language=Python, caption={Пример входной модели (фрагмент)}]
class WasteSiteVisitedContainers(BaseModel):
    plan = ForeignKey(WasteSiteVisitPlan, ...)
    container = ForeignKey(Container, ...)
    status = CharField(choices=ContainerStatus.choices)
    tour = ForeignKey(Tour, null=True, ...)
    is_moved = BooleanField(default=False)
    datetime_update = DateTimeField(...)
\end{lstlisting}

\subsection{Дополнительный контекст (опционально)}

\begin{itemize}
  \item Описание бизнес-правил в свободной форме
  \item Известные ограничения или инварианты
  \item Примеры проблем из прошлого опыта («раньше ловили баг когда..»)
\end{itemize}

Дополнительный контекст повышает качество анализа, но не обязателен.

% ════════════════════════════════════════════════════════════════
\section{Выходной отчёт}

\subsection{Структура}

Отчёт состоит из пяти категорий. Каждая находка содержит:
ссылку на узел графа, описание проблемы, и рекомендацию.

\begin{center}
\begin{tabularx}{\textwidth}{cLL}
  \toprule
  \textbf{Знак} & \textbf{Категория} & \textbf{Что сюда входит} \\
  \midrule
  \texttt{🔴} & Пробел компенсации
    & Ошибочная ветка не откатывает предыдущие шаги \\
  \texttt{❗} & Неожиданное состояние
    & Комбинация полей модели, которую граф не рассматривает \\
  \texttt{⚠️} & Крайний случай
    & Граничное входное значение или состояние БД \\
  \texttt{🟡} & Недостижимая ветка
    & Узел или ветка, недостижимые при нормальном потоке \\
  \texttt{ℹ️} & Замечание
    & Потенциальное улучшение без критичности \\
  \bottomrule
\end{tabularx}
\end{center}

\subsection{Пример отчёта}

\begin{warnbox}[title={🔴 Пробел компенсации — E20.1–E20.5}]
  \textbf{Узлы:} E20.1 (закрыть старые версии) → ERR1 (компенсация add\_tour\_task)\\
  \textbf{Проблема:} К моменту ERR1 шаги E20.1–E20.3 уже выполнены:
  старые версии контейнеров закрыты, новые созданы, зоны перенесены.
  Компенсация ERR1 описана как «уведомление об ошибке», но не включает
  откат версионирования контейнеров.\\
  \textbf{Рекомендация:} Добавить явный шаг отката E20.x в ветку ERR1,
  либо документировать что partial state допустим и обрабатывается отдельно.
\end{warnbox}

\begin{warnbox}[title={❗ Неожиданное состояние — E38.5}]
  \textbf{Узел:} E38.5 (проверки контейнера: статус плана, тип ТС, assignment, moved)\\
  \textbf{Проблема:} Поле \texttt{is\_moved=True} при \texttt{status=in\_route}
  — комбинация теоретически возможна если предыдущий вызов add\_tour завершился
  частично. Граф предполагает что \texttt{moved} и \texttt{in\_route} взаимоисключают,
  но инварианта в модели нет.\\
  \textbf{Рекомендация:} Добавить проверку или constraint в модель.
\end{warnbox}

\begin{infobox}[title={🟡 Недостижимая ветка — F2}]
  \textbf{Узел:} F1 → F2 (пересчитать агрегаты source\_tour\_ids)\\
  \textbf{Проблема:} Путь E25 (только пересортировать планы) ведёт в F0,
  после чего F1 спрашивает «были контейнеры на перенос?». При пути через E25
  контейнеров на перенос не было по определению (C3→Нет). Ветка F2 при этом
  пути никогда не активируется.\\
  \textbf{Рекомендация:} Уточнить нужен ли пересчёт агрегатов при пути E25→F0.
\end{infobox}

\subsection{Метрики отчёта}

В конце отчёта — сводная таблица:

\begin{center}
\begin{tabular}{lrr}
  \toprule
  \textbf{Категория} & \textbf{Найдено} & \textbf{Критичных} \\
  \midrule
  Пробелы компенсации & 2 & 2 \\
  Неожиданные состояния & 4 & 1 \\
  Крайние случаи & 6 & 2 \\
  Недостижимые ветки & 1 & 0 \\
  Замечания & 3 & 0 \\
  \midrule
  \textbf{Итого} & \textbf{16} & \textbf{5} \\
  \bottomrule
\end{tabular}
\end{center}

% ════════════════════════════════════════════════════════════════
\section{Архитектура}

\subsection{Место в системе}

Инструмент реализуется как поток \texttt{граф-аналитик} в
\texttt{development\_system} или как standalone Python-скрипт.
Оба варианта используют одинаковое ядро анализа.

\begin{center}
\begin{tabular}{lll}
  \toprule
  \textbf{Вариант} & \textbf{Запуск} & \textbf{Когда использовать} \\
  \midrule
  Поток & \texttt{/запустить граф-аналитик <файл>} & интерактивно через Telegram \\
  Скрипт & \texttt{python analyze.py graph.md models.py} & CI/CD, batch \\
  \bottomrule
\end{tabular}
\end{center}

\subsection{Компоненты}

\begin{enumerate}
  \item \textbf{Парсер графа} — извлекает из Mermaid узлы, рёбра, условия,
    терминальные и error-узлы. Строит внутренний граф (adjacency list).

  \item \textbf{Парсер моделей} — извлекает из Python-кода классы, поля,
    типы, choices, null/blank, FK-связи. Строит словарь {модель → поля}.

  \item \textbf{Аналитический движок (LLM)} — получает граф и модели в
    структурированном промпте, возвращает находки по категориям.

  \item \textbf{Форматтер отчёта} — преобразует JSON-ответ LLM в
    читаемый текст (Markdown, plain text, или PDF через LaTeX).
\end{enumerate}

\subsection{Алгоритм анализа}

LLM получает несколько специализированных запросов (не один большой):

\begin{center}
\begin{tabularx}{\textwidth}{cLL}
  \toprule
  \textbf{Шаг} & \textbf{Запрос} & \textbf{Что ищем} \\
  \midrule
  1 & Анализ error-веток
    & Что произошло до ERR-узла? Что нужно откатить? \\
  2 & Анализ моделей vs граф
    & Какие комбинации полей граф не рассматривает? \\
  3 & Анализ граничных входов
    & Что если containers=[], date в прошлом, ro\_zone=None? \\
  4 & Анализ достижимости
    & Есть ли узлы до которых нельзя добраться? \\
  5 & Синтез
    & Ранжирование находок по критичности \\
  \bottomrule
\end{tabularx}
\end{center}

Разбивка на шаги снижает вероятность галлюцинаций: каждый запрос
сфокусирован, граф не теряется в большом контексте.

\subsection{Промпт-стратегия}

\begin{reqbox}{Системный промпт (схема)}
Ты — анализатор бизнес-процессов. Тебе дан граф в формате Mermaid
и Django-модели участвующие в процессе.

Твоя задача: найти [КАТЕГОРИЯ]. Отвечай ТОЛЬКО в формате JSON:
\{
  "findings": [
    \{"node": "E20.1", "category": "compensation\_gap",
     "description": "...", "recommendation": "...",
     "severity": "critical|warning|info"\}
  ]
\}

Не выдумывай проблемы которых нет. Если всё в порядке — верни пустой список.
\end{reqbox}

\begin{infobox}[title={Почему отдельные запросы, а не один}]
Один большой промпт «найди все проблемы» даёт размытые результаты.
Пять специализированных запросов дают конкретные, проверяемые находки.
Стоимость: $\approx5 \times 3000$ токенов $\approx \$0.05$ на один анализ.
\end{infobox}

% ════════════════════════════════════════════════════════════════
\section{Интерфейс}

\subsection{Telegram-поток}

\begin{lstlisting}[language={}, numbers=none, caption={Пример диалога}]
> /запустить граф-аналитик

< Пришли Mermaid-граф (или файл .md):

> [прикрепляет add_tour_process_mermaid.md]

< Получен граф: 45 узлов, 5 фаз, 3 error-узла.
  Теперь пришли Django-модели (Python-код):

> [прикрепляет models.py или вставляет код]

< Анализирую... (5 запросов к LLM)

< 📊 Отчёт: add_tour_process
  🔴 Пробелы компенсации: 2 (критичных: 2)
  ❗ Неожиданные состояния: 4 (критичных: 1)
  ⚠️  Крайние случаи: 6 (критичных: 2)
  🟡 Недостижимые ветки: 1
  ℹ️  Замечания: 3
  Итого критичных: 5

  [детальный отчёт ниже]
\end{lstlisting}

\subsection{Команды}

\begin{center}
\begin{tabular}{lp{9cm}}
  \toprule
  \textbf{Команда} & \textbf{Описание} \\
  \midrule
  \texttt{граф-аналитик <граф>} & Только граф, без моделей (упрощённый анализ) \\
  \texttt{граф-аналитик <граф> <модели>} & Полный анализ \\
  \texttt{граф-аналитик --категория компенсация} & Только одна категория \\
  \texttt{граф-аналитик --pdf} & Вернуть отчёт как PDF \\
  \bottomrule
\end{tabular}
\end{center}

% ════════════════════════════════════════════════════════════════
\section{Реализация}

\subsection{Парсер Mermaid}

Граф парсится регулярными выражениями — полный Mermaid-парсер не нужен,
достаточно извлечь узлы и рёбра из \texttt{flowchart}:

\begin{lstlisting}[language=Python, caption={Парсер графа}]
import re
from dataclasses import dataclass, field
from typing import Optional

@dataclass
class Node:
    id: str
    label: str
    is_error: bool = False
    is_terminal: bool = False
    is_decision: bool = False

@dataclass
class Edge:
    from_id: str
    to_id: str
    condition: Optional[str] = None

def parse_mermaid(text: str) -> tuple[dict[str, Node], list[Edge]]:
    nodes, edges = {}, []

    # Узлы: A0["E1 Текст"] или A0{"Вопрос?"} или A0(["Текст"])
    for m in re.finditer(r'(\w+)\[\"([^\"]+)\"\]', text):
        nid, label = m.group(1), m.group(2)
        nodes[nid] = Node(
            id=nid, label=label,
            is_error='ERR' in nid or 'ошибк' in label.lower(),
            is_terminal='финиш' in label.lower() or 'completed' in label.lower(),
        )
    for m in re.finditer(r'(\w+)\{\"([^\"]+)\"\}', text):
        nid, label = m.group(1), m.group(2)
        nodes[nid] = Node(id=nid, label=label, is_decision=True)

    # Рёбра: A --> B или A -->|"условие"| B
    for m in re.finditer(r'(\w+)\s*-->(?:\|"([^\"]+)"\|)?\s*(\w+)', text):
        edges.append(Edge(m.group(1), m.group(3), m.group(2)))

    return nodes, edges
\end{lstlisting}

\subsection{Парсер моделей}

\begin{lstlisting}[language=Python, caption={Парсер Django-моделей}]
import ast

def parse_models(source: str) -> dict[str, list[dict]]:
    """Извлекает {ClassName: [{name, type, kwargs}]} из Python-кода."""
    tree = ast.parse(source)
    result = {}
    for node in ast.walk(tree):
        if not isinstance(node, ast.ClassDef):
            continue
        fields = []
        for stmt in node.body:
            if not isinstance(stmt, ast.Assign):
                continue
            for target in stmt.targets:
                if not isinstance(target, ast.Name):
                    continue
                fields.append({
                    'name': target.id,
                    'type': ast.unparse(stmt.value)
                })
        if fields:
            result[node.name] = fields
    return result
\end{lstlisting}

\subsection{Аналитический движок}

\begin{lstlisting}[language=Python, caption={Запрос к LLM (один шаг)}]
import json, httpx, os

PROMPTS = {
    'compensation': """
Граф процесса (Mermaid):
{graph}

Django-модели:
{models}

Найди пробелы в компенсации ошибок: случаи когда error-ветка
(ERR-узлы) не откатывает шаги уже выполненные до неё.
Отвечай JSON: {{"findings": [{{"node":"...", "description":"...",
"recommendation":"...", "severity":"critical|warning"}}]}}
""",
    'edge_cases': """...""",
    'model_states': """...""",
    'reachability': """...""",
}

def analyze_step(step: str, graph: str, models: str) -> list[dict]:
    prompt = PROMPTS[step].format(graph=graph, models=models)
    resp = httpx.post(
        "https://openrouter.ai/api/v1/chat/completions",
        headers={"Authorization": f"Bearer {os.environ['OPENROUTER_API_KEY']}"},
        json={"model": "anthropic/claude-opus-4-5",
              "messages": [{"role": "user", "content": prompt}]},
        timeout=60,
    ).json()
    content = resp['choices'][0]['message']['content']
    # Извлечь JSON из ответа
    m = re.search(r'\{.*\}', content, re.DOTALL)
    return json.loads(m.group())['findings'] if m else []

def analyze(graph_text: str, models_text: str) -> dict:
    all_findings = []
    for step in PROMPTS:
        all_findings.extend(analyze_step(step, graph_text, models_text))
    return {
        'total': len(all_findings),
        'critical': sum(1 for f in all_findings if f.get('severity') == 'critical'),
        'findings': all_findings,
    }
\end{lstlisting}

\subsection{Структура файлов}

\begin{lstlisting}[language={}, numbers=none]
development_system/
  мастер/потоки/
    граф-аналитик.lisp     ← поток (обёртка над Python)
  tools/
    graph_analyzer/
      __init__.py
      parser.py             ← parse_mermaid, parse_models
      analyzer.py           ← analyze(), analyze_step()
      prompts.py            ← все промпты
      formatter.py          ← форматирование отчёта
      cli.py                ← python -m graph_analyzer <файлы>
\end{lstlisting}

% ════════════════════════════════════════════════════════════════
\section{Критерии приёмки}

\subsection{Функциональные}

\begin{enumerate}
  \item Принимает граф add\_tour (45 узлов) и находит
    \textbf{не менее 3 реальных проблем}, подтверждённых разработчиком
  \item Каждая находка содержит ссылку на конкретный узел графа
  \item Отчёт возвращается за \textbf{не более 60 секунд}
  \item Работает без дополнительного контекста (только граф + модели)
  \item Ложноположительные находки \textbf{не более 30\%} от общего числа
\end{enumerate}

\subsection{Нефункциональные}

\begin{enumerate}
  \item Стоимость одного анализа \textbf{не более \$0.10} (OpenRouter)
  \item Запуск из Telegram одной командой
  \item Результат читается без технических знаний о промптах
\end{enumerate}

\subsection{Не входит в v1}

\begin{itemize}
  \item Генерация тестов (этап 2, отдельное ТЗ)
  \item Анализ кода реализации (только граф и модели)
  \item Сохранение истории анализов
  \item Сравнение двух версий графа
\end{itemize}

% ════════════════════════════════════════════════════════════════
\section{Риски}

\begin{center}
\begin{tabularx}{\textwidth}{lLl}
  \toprule
  \textbf{Риск} & \textbf{Описание} & \textbf{Митигация} \\
  \midrule
  Галлюцинации LLM
    & Находки которых нет в реальности
    & Разработчик проверяет каждую; порог ЛП \textless 30\% \\
  Большой граф
    & 45+ узлов перегружают контекст
    & Анализ пофазно: каждая subgraph отдельно \\
  Неполные модели
    & Пользователь не передал все поля
    & Инструмент явно сообщает каких полей не хватает \\
  Стоимость
    & Частые запуски дорожают
    & Кэш результатов по хешу входа \\
  \bottomrule
\end{tabularx}
\end{center}

% ════════════════════════════════════════════════════════════════
\section{Открытые вопросы}

\begin{enumerate}
  \item \textbf{Модель LLM:} claude-opus-4-5 даёт лучшее качество,
    gpt-4.1 — дешевле. Решить по результатам первого прогона на add\_tour.
  \item \textbf{Язык реализации:} Python (проще интеграция с Django),
    или CL-поток (нативен для development\_system). Рекомендация: Python-скрипт
    + тонкая CL-обёртка для вызова из Telegram.
  \item \textbf{Формат отчёта:} plain text в Telegram достаточен для v1,
    PDF — для v2.
  \item \textbf{Фазный анализ:} анализировать граф целиком или пофазно
    (по subgraph)? Пофазно дешевле и точнее — но требует маппинга фаз.
\end{enumerate}

% ════════════════════════════════════════════════════════════════
\section{Следующий шаг}

\begin{okbox}[title={MVP за 1 день}]
\begin{enumerate}
  \item Python-скрипт \texttt{analyze.py}: парсер Mermaid + 2 промпта
    (компенсация + крайние случаи)
  \item Прогнать на add\_tour + WasteSiteVisitedContainers
  \item Сравнить находки с реальными проблемами известными команде
  \item Если точность \textgreater 50\% — добавить остальные 3 промпта
\end{enumerate}
\end{okbox}

\vfill
\hrule
\vspace{4mm}
{\small\color{MidGray}
  ТЗ v1.0 · 14~марта~2026 · development\_system ·
  Следующий документ: ТЗ на генератор тестов (этап~2)
}

\end{document}
